\SourceLeaf{248}
\Article{373}
\textsc{Problema.} \emph{Si functio $X$ secundum modulum $p$ in factores inter se primos $\xi,\xi',\xi''$ etc. sit resoluta, $X$ secundum modulum $pp$ in similes factores $\Xi,\Xi',\Xi''$ etc. resolvere ita, ut sit}
\[
\xi\equiv\Xi,\quad\xi'\equiv\Xi',\quad\xi''\equiv\Xi'',\text{ etc.}\;(\mathrm{mod.}\,p)
\]

\emph{Sol.} Sit $X\equiv\xi\psi(\mathrm{mod.}\,p)$ seu $X=\xi\psi+p\Sigma$. Ponatur
\[
\Xi=\xi+p\varphi,\quad\Psi=\psi+p\omega
\]
erit
\[
\Xi\Psi=X-p\Sigma+(\varphi\psi+\xi\omega)p+pp\varphi\omega
\]
Si igitur $\Xi\Psi$ esse debet $\equiv X(\mathrm{mod.}\,pp)$, necessario debet esse $\varphi\psi+\xi\omega-\Sigma$ per $p$ divisibilis. At cum $\psi$ et $\xi$ secundum modulum $p$ sint functiones inter se primae, $\varphi$ et $\omega$ ita determinari poterunt, ut haec conditio adimpleatur (\textsection.\ 336), et quidem insuper ita, ut dimensiones ipsarum $\varphi$ et $\omega$ sint respective unitate minores dimensionibus functionum $\xi,\psi$. Hinc erit $X\equiv\Xi\Psi(\mathrm{mod.}\,pp)$. Patet, simili modo $\Psi$ rursus in factores $\Xi'\Omega$ discerpi posse, ita ut alter $\Xi'$ sit $\equiv\xi'(\mathrm{mod.}\,p)$ et ita porro, unde tandem
\[
X\equiv\Xi\,\Xi'\,\Xi''\text{ etc.}\;(\mathrm{mod.}\,pp).\quad\text{Q. E. Fac.}
\]

\Article{374}
Facile hinc probari potest, functionem $X$ etiam secundum modulos $p^3,p^4$ etc. in factores resolvi posse. Generaliter sit
\[
X\equiv PQ(\mathrm{mod.}\,p^m)\quad\text{seu}\quad X=PQ+p^mR
\]
et functio $P$ ad ipsam $Q$ prima secundum modulum $p$; posito
\[
P'=P+Ap^m,\quad Q'=Q+Bp^m
\]
erit
\[
P'Q'=X-p^mR+(AQ+BP)p^m+ABp^{2m}
\]
Hinc pro quovis modulo $p^\nu$ ($\nu$ existente $>m$ et $<2m+1$) erit
\[
P'Q'\equiv X,\quad\text{si}\quad R\equiv AQ+BP(\mathrm{mod.}\,p^{\nu-m})
\]

\SourceLeaf{249}
Ex his perspicitur, si functio $X$ aequales non habeat divisores secundum modulum $p$, eam secundum modulum $p^k$ similiter in factores discerpi posse, uti secundum modulum $p$. At si $X$ divisores aequales habeat, res fit multo magis complicata neque adeo ex principiis praecedentibus prorsus exhauriri potest. Quare quum quae huc pertineant cuncta communicare non possimus, unicum casum tantummodo considerabimus, qui plurimum occurrit cuiusque enodatio ad quaedam in praecedentibus dubia solvenda requiritur. Hic est, si factores aequales unius dimensionis tantum respiciantur. Hic proprie etiam ad congruentiarum radices inveniendas adhiberi potest. Generaliter alia occasione hanc rem pertractabimus.

\Article{375}
Sit igitur $X\equiv X'(x-a)^m(\mathrm{mod.}\,p)$ et functio $X'$ ad $x-a$ prima: desiderantur omnes divisores unius dimensionis huic $x-a$ secundum modulum $p$ congrui ipsius $X$ secundum modulos $pp,p^3$ etc. (Supponimus, functionem $X$ absolute per $x-a$ dividi non posse; alias enim $x-a$ secundum modulum quemcunque functionem $X$ divideret). Si substituatur $z+a$ pro $x$, habebitur
\[
Z\equiv Z'z^m(\mathrm{mod.}\,p)\quad\text{seu}\quad Z=Z'z^m+pA
\]
Iam si $Z$ secundum modulum $pp$ per aliquem divisorem formae $z+\alpha p$ dividi potest, necessario $A$ debet esse formae $zZ''+pB$. Nisi hoc sit, disquisitio iam est finita. Ponamus igitur
\[
\begin{aligned}
Z&\equiv Z'z^m+pZ''z(\mathrm{mod.}\,pp)\\
&\text{seu}\quad Z=Z'z^m+pZ''z+ppB
\end{aligned}
\]
patetque, $Z$ per $z$ ac quemcunque alium divisorem huic secundum modulum $p$ congruum dividi posse.

Ut attentio fixetur, ponemus $m=4$, facile perspicitur, quemvis alium casum simili modo tractari posse. Iam si $Z$ secundum modulum $p^3$ per aliquem divisorem formae $z+\alpha p$ dividi potest, erit
\[
\begin{aligned}
0&\equiv-\alpha ppZ''+ppB(\mathrm{modd.}\,z+\alpha p,p^3)\\
&\text{seu}\quad\alpha Z''\equiv B(\mathrm{modd.}\,z,p)
\end{aligned}
\]
Iam tres casus esse possunt

1) si $Z''\equiv0(\mathrm{modd.}\,z,p)$ et $B$ non $\equiv0$, tunc patet, nullum ipsius $\alpha$ valorem congruentiae satisfacere adeoque $Z$ secundum modulum $p^3$ nullum divisorem formae $z+\alpha p$ habere. Quare disquisitio erit finita

\SourceLeaf{250}
2) si nec $Z''$ nec $B\equiv0(\mathrm{modd.}\,z,p)$; tunc $\alpha$ unicum valorem habebit, scilicet
\[
\alpha\equiv\frac B{Z''}(\mathrm{modd.}\,z,p).
\]
Quare erit unicus divisor $\equiv z+\alpha p(\mathrm{mod.}\,pp)$ ipsius $Z$ secundum modulum $p^3$; eritque
\[
Z\equiv V(z+\alpha p)+p^3W
\]
Iam ponatur divisor ipsius $Z(\mathrm{mod.}\,p^4)$ $z+\alpha p+\beta pp$ eritque
\[
0\equiv\qquad\qquad
\]
\BookRule{.28\linewidth}
